% cas-sc nạp gói STIX kiểu Type 1, dễ sai ánh xạ ký hiệu khi biên dịch XeLaTeX.
% Đánh dấu gói đã nạp để dùng bộ math mặc định ổn định hơn với fontspec.
\makeatletter
\@namedef{ver@stix.sty}{9999/12/31 disabled for XeLaTeX}
\makeatother
\documentclass[a4paper,fleqn]{cas-sc}

\usepackage[vietnamese,english]{babel}
\usepackage{fontspec}
% Gọi theo tên file để cả Tectonic lẫn Overleaf tìm font trong TeX Live.
\setmainfont{texgyretermes}[
  Extension=.otf,
  UprightFont=*-regular,
  BoldFont=*-bold,
  ItalicFont=*-italic,
  BoldItalicFont=*-bolditalic]
\setsansfont{texgyreheros}[
  Extension=.otf,
  UprightFont=*-regular,
  BoldFont=*-bold,
  ItalicFont=*-italic,
  BoldItalicFont=*-bolditalic]
\setmonofont{texgyrecursor}[
  Extension=.otf,
  UprightFont=*-regular,
  BoldFont=*-bold,
  ItalicFont=*-italic,
  BoldItalicFont=*-bolditalic]
\usepackage{amsmath}
\usepackage{booktabs}
\usepackage{graphicx}
\usepackage[numbers]{natbib}
\usepackage{xurl}

% Cho phép biên dịch trực tiếp report/main.tex hoặc qua report_template.tex.
\providecommand{\ReportPrefix}{}
\graphicspath{{\ReportPrefix figures/}}
\input{\ReportPrefix generated/environment.tex}

\newcommand{\StudentName}{Bùi Đình Cảnh}
\newcommand{\StudentID}{24021392}
\newcommand{\ClassCode}{INT3011E\_3}
\newcommand{\StudentEmail}{buidinhcanh2006@gmail.com}
\newcommand{\StudentAffiliation}{Sinh viên \StudentID, lớp \ClassCode, Trường Đại học Công nghệ, Đại học Quốc gia Hà Nội; email: \StudentEmail}
\newcommand{\RepositoryURL}{https://github.com/WilliamsBui-fnd/INT3011E_3}

\title[mode=title]{Thực nghiệm giải N-Queens bằng SAT, CP và MIP đến N=1000}
\shorttitle{Giải N-Queens bằng SAT, CP và MIP}

\author[1]{\StudentName}
\shortauthors{\StudentName}
\ead{\StudentEmail}
\address[1]{\StudentAffiliation}

\begin{document}

\begin{abstract}
Bài báo nghiên cứu các mô hình giải chính xác cho bài toán N-Queens, với trọng
tâm là ảnh hưởng của kỹ thuật mã hóa ràng buộc at-most-one trong SAT. Năm mã
hóa Pairwise, Sequential Counter, Commander, Product và Binary được triển khai
trên PySAT/Glucose3 và so sánh với OR-Tools CP-SAT, Gurobi MIP, CPLEX MIP và
CPLEX CP Optimizer. Quy trình thực nghiệm có timeout cưỡng bức, checkpoint,
kiểm tra nghiệm và phân biệt timeout, giới hạn giấy phép và tiến trình bị hệ
điều hành dừng. Bộ dữ liệu sau khi bổ sung Gurobi Academic WLS gồm 241 lượt
chạy trên các kích thước từ $N=1$ đến $N=1000$. Tất cả 182 nghiệm SAT thu được
đều vượt qua bộ kiểm tra độc lập; 14
lượt UNSAT tại $N=2,3$ cũng phù hợp với tính chất bài toán. Trên môi trường
Google Colab đã sử dụng, Gurobi hoàn thành 5/5 lượt tại $N=64,128$ và 3/3 lượt
tại $N=256$. Product là mã hóa SAT duy nhất hoàn thành cả ba lần tại $N=1000$,
với trung vị 28.127 giây, 1,298,508 biến và 11,942,164 mệnh đề.
Kết quả cũng cho thấy kích thước CNF không đủ để dự đoán thời gian giải: hiệu
năng thay đổi không đơn điệu theo $N$ và phụ thuộc mạnh vào cấu trúc encoding
cũng như heuristic của solver.
\end{abstract}

\begin{keywords}
N-Queens \sep SAT solving \sep cardinality encoding \sep PySAT \sep
constraint programming \sep mixed-integer programming
\end{keywords}

\maketitle

\section{Giới thiệu}

Bài toán N-Queens yêu cầu đặt $N$ quân hậu trên bàn cờ có $N$ hàng và $N$ cột sao cho
không có hai quân nào cùng hàng, cùng cột hoặc cùng đường chéo. Bài toán có mô
tả ngắn nhưng là một ví dụ điển hình để khảo sát ảnh hưởng của mô hình hóa tới
hiệu năng solver. Cùng một tập ràng buộc có thể được biểu diễn bằng công thức
SAT ở dạng CNF, bằng các ràng buộc toàn cục trong Constraint Programming (CP),
hoặc bằng hệ ràng buộc tuyến tính trên biến nhị phân trong MIP.

Sự khác biệt giữa ba biểu diễn nằm ở cấu trúc mà solver có thể khai thác.
SAT yêu cầu chuyển ràng buộc đếm thành CNF và có thể tạo nhiều biến phụ; CP
giữ được ràng buộc \texttt{AllDifferent}; MIP dùng biến nhị phân cùng các bất
đẳng thức tuyến tính. Bài toán vì thế phù hợp để khảo sát chi phí của mô hình
hóa tách biệt với việc tìm một nghiệm khả thi.

Mục tiêu của nghiên cứu không phải đếm mọi lời giải, mà tìm một cấu hình hợp lệ
hoặc chứng minh không tồn tại cấu hình. Đóng góp chính gồm: (i) triển khai và
kiểm thử năm SAT encoding cho AMO; (ii) xây dựng bốn mô hình đối chứng bằng
OR-Tools, Gurobi và CPLEX; (iii) xây dựng pipeline thực nghiệm có thể tiếp tục
sau khi Colab bị ngắt; và (iv) đánh giá đồng thời thời gian, quy mô encoding và
tỷ lệ hoàn thành đến $N=1000$.

\section{Bối cảnh và câu hỏi nghiên cứu}

Các solver SAT hiện đại kết hợp tìm kiếm CDCL và học mệnh đề; hiệu năng của
chúng phụ thuộc vào công thức CNF được cung cấp~\cite{glucose2009}. Đối với
ràng buộc at-most-one (AMO), Pairwise ngắn gọn và không dùng biến phụ, trong
khi Sequential Counter dùng biến trạng thái để giảm số clause
~\cite{sinz2005}. PySAT cung cấp nhiều kiểu cardinality encoding với cùng giao
diện~\cite{pysat-paper,pysat-docs}. N-Queens cũng có mô hình CP tự nhiên với
ba ràng buộc \texttt{AllDifferent}~\cite{ortools-nqueens}.

Báo cáo trả lời ba câu hỏi: (i) kích thước CNF của năm cách mã hóa AMO thay
đổi như thế nào theo $N$; (ii) sự thay đổi đó có dự báo được thời gian giải
hay không; và (iii) SAT, CP-SAT và MIP khác nhau ra sao về tỷ lệ hoàn thành
trong cùng giới hạn thời gian. So sánh thực nghiệm được giới hạn bởi những
giá trị $N$ đã chạy được và điều kiện cấp phép của từng solver.

\section{Mô hình bài toán}

Với mỗi ô $(r,c)$, đặt biến Boolean $x_{r,c}$ bằng 1 khi có hậu tại ô đó. Mô
hình SAT và MIP sử dụng các nhóm ràng buộc sau:
\begin{align}
  \sum_{c=0}^{N-1} x_{r,c} &= 1 && \text{với mọi hàng }r, \\
  \sum_{r=0}^{N-1} x_{r,c} &= 1 && \text{với mọi cột }c, \\
  \sum_{(r,c)\text{ thuộc }d} x_{r,c} &\leq 1
    && \text{với mọi đường chéo }d.
\end{align}

Mỗi ô có đúng một biến cơ sở, nên mô hình SAT/MIP bắt đầu với $N^2$ biến.
Mô hình MIP cài đặt $2N$ ràng buộc hàng/cột và $4N-2$ ràng buộc đường chéo,
trong đó các đường chéo chỉ có một ô tạo bất đẳng thức dư nhưng không làm
thay đổi tập nghiệm. SAT chỉ sinh AMO cho đường chéo dài ít nhất hai ô.

Trong SAT, ràng buộc bằng 1 được tách thành at-least-one (ALO) và at-most-one
(AMO). Với nhóm literal $X=(x_1,\ldots,x_k)$, ALO là clause
$(x_1\mathbin{\mathrm{OR}}\cdots\mathbin{\mathrm{OR}}x_k)$. AMO được thay đổi
theo từng encoding. Trong CP, mỗi hàng $i$ dùng một biến nguyên $q_i$ nhận giá
trị từ 0 đến $N-1$ để biểu diễn cột của
quân hậu; mô hình chỉ cần ba ràng buộc toàn cục
\texttt{AllDifferent} trên $q_i$, $q_i+i$ và $q_i-i$, đúng với mô hình N-Queens
được OR-Tools công bố~\cite{ortools-nqueens}.

Hai mô hình biểu diễn cùng một tập cấu hình: hàng $r$ được gán cột $q_r=c$
đúng khi $x_{r,c}=1$. \texttt{AllDifferent} trên $q_r$ loại xung đột cột;
hai biểu thức còn lại lần lượt loại hai hướng đường chéo. Điều này cho phép
kiểm tra nghiệm do mọi solver trả về bằng cùng một hàm xác minh độc lập.

\section{Các phương pháp giải}

\subsection{SAT và các mã hóa AMO}

PySAT cung cấp giao diện thống nhất cho solver và các encoding cardinality
\cite{pysat-docs,pysat-paper}. Glucose3 được dùng làm SAT backend. Pairwise
sinh clause $(\overline{x_i}\mathbin{\mathrm{OR}}\overline{x_j})$ cho mọi cặp. Sequential Counter dùng
biến trạng thái phụ tuyến tính theo kích thước nhóm, dựa trên ý tưởng của
Sinz~\cite{sinz2005}. Binary gán mã bit cho vị trí của literal đúng. Commander
chia các literal thành nhóm và tạo biến ``chỉ huy'' cho mỗi nhóm. Product ánh
xạ literal lên một lưới phụ hai chiều và áp dụng AMO cho các chỉ báo hàng/cột.

Pairwise, Sequential và Binary gọi trực tiếp \texttt{CardEnc} của PySAT;
Commander và Product được cài đặt trong project. ALO được áp dụng cho mọi hàng
và cột; AMO được áp dụng cho hàng, cột và đường chéo. Bộ kiểm thử truth-table
kiểm tra trực tiếp rằng cả năm cài đặt AMO chấp nhận đúng các phép gán có không
quá một literal đúng. Ngoài ra, Glucose3 là solver CDCL có heuristic học clause;
những đặc trưng này giải thích vì sao số clause không ánh xạ đơn giản thành thời
gian chạy~\cite{glucose2009}.

Với nhóm $X$ gồm $k$ literal, Pairwise tạo $\binom{k}{2}$ clause nhị phân và
không cần biến phụ. Sequential Counter biểu diễn thông tin tiền tố đã có
literal đúng bằng biến phụ. Binary gán cho literal đúng một mã dài
$\lceil\log_2 k\rceil$ bit. Bảng~\ref{tab:amo-complexity} tóm tắt bậc kích
thước của từng nhóm AMO theo đúng cách cài đặt trong project; ALO bổ sung đúng
một clause độ dài $k$ cho mỗi hàng/cột.

\begin{table}[htbp]
  \centering
  \caption{Bậc kích thước của một nhóm AMO gồm $k$ literal. Các hằng số và
  trường hợp nhóm rất nhỏ phụ thuộc cài đặt cụ thể.}
  \label{tab:amo-complexity}
  \begin{tabular}{lcc}
    \toprule
    Mã hóa & Biến phụ & Clause \\
    \midrule
    Pairwise & $0$ & $O(k^2)$ \\
    Sequential & $O(k)$ & $O(k)$ \\
    Binary & $O(\log k)$ & $O(k\log k)$ \\
    Commander & $O(\sqrt{k})$ & $O(k\sqrt{k})$ \\
    Product & $O(\sqrt{k})$ & $O(k)$ \\
    \bottomrule
  \end{tabular}
\end{table}

Trong Commander, $k$ literal được chia thành các nhóm kích thước xấp xỉ
$\sqrt{k}$. Mỗi nhóm có một biến chỉ huy $c_g$; hai chiều kéo theo giữa
$c_g$ và phép OR của nhóm bảo đảm $c_g$ đúng chính xác khi nhóm có literal
đúng. AMO Pairwise được áp dụng bên trong mỗi nhóm và giữa các chỉ huy. Product
đặt các literal lên lưới phụ $p\times q$ với $p=\lceil\sqrt{k}\rceil$,
$q=\lceil k/p\rceil$. Một literal đúng kéo theo hai biến chỉ báo hàng và cột
tương ứng; AMO Pairwise trên hai tập chỉ báo ngăn đồng thời hai literal đúng.
Với nhóm tối đa bốn literal, bản cài đặt dùng Pairwise trực tiếp.

\subsection{Lập luận bảo toàn nghiệm}

ALO cùng AMO trên mỗi hàng và cột buộc đúng một hậu tại mỗi hàng/cột. AMO
trên hai họ đường chéo loại mọi cặp hậu tấn công nhau. Vì vậy phép gán CNF
thỏa mãn luôn cho một bàn cờ hợp lệ. Chiều ngược lại, từ một bàn cờ hợp lệ
đặt các biến ô tương ứng bằng đúng và các biến ô còn lại bằng sai; mọi nhóm
ALO/AMO đều thỏa. Các biến phụ của Sequential và Binary có thể được gán
theo encoding của PySAT. Với Commander, đặt chỉ huy đúng cho nhóm chứa
literal đúng. Với Product, đặt đúng chỉ báo hàng/cột chứa literal đúng. Do
đó mọi bàn cờ hợp lệ đều mở rộng thành một phép gán CNF thỏa mãn.

Kiểm thử truth-table xét toàn bộ $2^5$ phép gán của nhóm năm literal cho cả
năm encoding và so sánh kết quả solver với điều kiện nhiều nhất một bit
đúng. Kiểm thử đầu-cuối tiếp tục xác nhận SAT cho $N=1,4,\ldots,8$ và
UNSAT cho $N=2,3$, đồng thời kiểm tra mọi nghiệm bằng hàm độc lập.

\subsection{CP-SAT và CP Optimizer}

Mô hình OR-Tools CP-SAT và CPLEX CP Optimizer đều dùng $N$ biến nguyên cùng ba
ràng buộc \texttt{AllDifferent}. CP-SAT là solver CP trên nền SAT cho biến
nguyên~\cite{ortools-cpsat}; CP Optimizer cung cấp các ràng buộc toàn cục cho
bài toán tổ hợp~\cite{ibm-cp}. Mỗi solver được giới hạn một worker để giảm sai
lệch do mức độ song song khác nhau.

Mô hình CP không cần liệt kê $N^2$ biến ô. Mỗi $q_r\in\{0,\ldots,N-1\}$
chọn cột cho hàng $r$; ba ràng buộc \texttt{AllDifferent} lần lượt áp dụng
cho các dãy $(q_r)$, $(q_r+r)$ và $(q_r-r)$. Cùng một biểu diễn mức cao
được đưa cho hai engine CP, song cơ chế lan truyền và tìm kiếm nội bộ khác
nhau nên thời gian vẫn có thể khác đáng kể.

\subsection{Gurobi MIP và CPLEX MIP}

Hai mô hình MIP dùng $N^2$ biến nhị phân và các phương trình/bất phương trình
ở phần trước. Đây là bài toán khả thi, vì vậy không cần hàm mục tiêu. Gurobi và
CPLEX được cấu hình một thread, seed theo lần lặp và time limit nội bộ. Gurobi
tuân theo mô hình formulate--solve--query của API tối ưu toán học
\cite{gurobi-manual}; CPLEX MIP dùng cùng mô hình để tạo phép so sánh giữa hai
solver thương mại.

Sau khi solver dừng, nếu có nghiệm khả thi, chương trình lấy các ô có giá trị
lớn hơn $0{,}5$ và kiểm tra lại số hậu, miền tọa độ, hàng, cột và đường chéo.
Trạng thái hết giờ được phân biệt với vô nghiệm đã chứng minh; lỗi license
được ghi thành \texttt{LICENSE\_LIMIT}, không gán thời gian giải bằng không.

\section{Thiết kế thực nghiệm}

Thực nghiệm chạy trên Google Colab với \ResultCpuCount{} CPU logic, một worker
cho mỗi solver và bốn phase:
\begin{itemize}
  \item smoke: $N=1,2,3,4,8$, một lần, timeout 15 giây;
  \item standard: $N=16,32,64,128$, năm lần, timeout 60 giây;
  \item CPLEX: $N=8,16,32$, ba lần, timeout 60 giây;
  \item large: $N=256,512,1000$, ba lần, timeout 120 giây.
\end{itemize}

Môi trường ghi nhận Linux, Python \ResultPython{}, PySAT \ResultPySAT{},
OR-Tools \ResultORTools{}, Gurobi \ResultGurobi{} và CPLEX \ResultCPLEX{}.
Pairwise không được chạy ở phase large vì số clause tăng mạnh. Sau đợt đo
ban đầu, giấy phép Gurobi Academic WLS đã cho phép chạy lại $N=64,128$,
năm lần mỗi kích thước, và chạy bổ sung $N=256$ ba lần. Các file kết quả
ban đầu được giữ nguyên; khi gộp, mỗi run Gurobi mới thay thế đúng khóa
$(phase,N,method,run)$ cũ bị \texttt{LICENSE\_LIMIT}. Vì vậy bộ gộp có
thêm ba run ở $N=256$, không nhân đôi mười run tại $N=64,128$.

\texttt{Total\_Time} gồm thời gian xây dựng mô hình/encoding và thời gian giải,
nhưng không gồm chi phí khởi tạo process cha. \texttt{Wall\_Time} bao gồm cả
khởi tạo process và IPC. Với Glucose3, seed trong script không được truyền vào
solver; vì vậy các lần lặp SAT đo độ dao động thời gian trên cùng cấu hình tìm
kiếm mặc định, không phải các lần chạy ngẫu nhiên độc lập.

Mỗi lượt chạy được thực hiện trong process riêng. Process cha đợi tối đa
giới hạn thời gian đã chọn, sau đó dừng process con nếu cần; ngay sau mỗi
lượt, một dòng CSV được ghi và đồng bộ xuống đĩa để có thể tiếp tục khi Colab
ngắt. Chỉ các run có trạng thái SAT hoặc UNSAT mới được dùng để tính trung vị,
min và max của \texttt{Total\_Time}; tỷ lệ thành công được tính trên toàn bộ
run, kể cả timeout và lỗi. Không nội suy thời gian cho run chưa hoàn thành.

Các run bổ sung dùng một phiên Colab khác đợt đo gốc. Metadata mới ghi nhận
2 CPU logic và 12{,}67 GiB RAM, nhưng metadata cũ không lưu model CPU và
tổng RAM. Vì vậy các so sánh thời gian giữa Gurobi bổ sung và các solver đã
đo trước đó chỉ mang tính mô tả, không phải một phép so sánh phần cứng được
kiểm soát chặt.

\section{Kết quả và thảo luận}

\subsection{Tính đúng đắn và mức độ hoàn thành}

Bộ dữ liệu có 241 run, không có khóa $(phase,N,method,run)$ bị trùng. Tổng hợp
trạng thái được trình bày trong Bảng~\ref{tab:status}. Tất cả 182 nghiệm SAT
đều có \texttt{Solution\_Valid=True}; 14 kết quả UNSAT đều thuộc $N=2$ và
$N=3$. Không có \texttt{INVALID\_SOLUTION}. Ba run
\texttt{LICENSE\_LIMIT} còn lại đều thuộc CPLEX MIP, không thuộc Gurobi.

\begin{table}[htbp]
  \centering
  \caption{Phân bố trạng thái của toàn bộ lượt chạy.}
  \label{tab:status}
  \input{\ReportPrefix generated/status_table.tex}
\end{table}

Hình~\ref{fig:status} thể hiện số lần thành công trên tổng số lần chạy. Cần đọc
hình này cùng đồ thị thời gian: trung vị chỉ được tính trên các run thành công,
nên một điểm thời gian không biểu diễn xác suất hoàn thành.

\begin{figure*}[htbp]
  \centering
  \includegraphics[width=\textwidth]{plot_status.png}
  \caption{Ma trận hoàn thành. TO: timeout; LIC: giới hạn license; KILL:
  tiến trình nhận SIGKILL. Ô trống là cấu hình không được đưa vào phase đó.}
  \label{fig:status}
\end{figure*}

\subsection{Thời gian giải}

Hình~\ref{fig:time} và Bảng~\ref{tab:fastest} cho thấy không có một phương pháp
thắng ở mọi kích thước. Sequential có trung vị thấp nhất tại $N=16$; Gurobi
MIP tại $N=32,64,128$; Product tại $N=256,1000$; Binary tại $N=512$.
Với $N=64,128$, Gurobi hoàn thành 5/5 run, trung vị lần lượt 0{,}098 và
0{,}194 giây. Tại $N=256$, Gurobi hoàn thành 3/3 run với trung vị
3{,}078 giây, còn Product là 2{,}704 giây. Những thứ hạng này mô tả bộ dữ
liệu đã gộp, không khẳng định ưu thế tổng quát giữa các solver.

\begin{figure*}[htbp]
  \centering
  \includegraphics[width=\textwidth]{plot_time.png}
  \caption{Trung vị thời gian của các run thành công; whisker là min--max.
  Đường không được nối qua kích thước không có run thành công.}
  \label{fig:time}
\end{figure*}

\begin{table}[htbp]
  \centering
  \caption{Phương pháp có trung vị thấp nhất trong các cấu hình đã chạy.}
  \label{tab:fastest}
  \resizebox{\linewidth}{!}{\input{\ReportPrefix generated/fastest_table.tex}}
\end{table}

Bảng~\ref{tab:gurobi-comparison} đặt các run mới cạnh những phương pháp
đã chạy trên cùng kích thước. Tại $N=64$, Product và Binary đều timeout
5/5 lần trong khi Pairwise, CP-SAT và Gurobi hoàn thành. Ở $N=128$ cả
Product và Binary đều trở lại thành công 5/5 lần, nhưng trung vị vẫn lớn
hơn Gurobi trong bộ dữ liệu này. Tại $N=256$, CP-SAT chỉ hoàn thành 2/3
run và Pairwise không được đưa vào phase lớn.

\begin{table*}[htbp]
  \centering
  \caption{Trung vị thời gian theo giây của các run thành công; trong ngoặc
  là số run thành công/tổng số run. TO: toàn bộ run timeout; dấu gạch ngang:
  không chạy.}
  \label{tab:gurobi-comparison}
  \resizebox{\textwidth}{!}{\input{\ReportPrefix generated/gurobi_comparison_table.tex}}
\end{table*}

Bảng~\ref{tab:sat-time} tập trung vào SAT. Product timeout ở $N=32,64$ nhưng
hoàn thành cả ba lần ở $N=1000$; Binary timeout tại $N=64$ nhưng hoàn thành tại
$N=128,256,512$. Đây là bằng chứng thực nghiệm rằng độ khó tìm kiếm không tăng
đơn điệu theo $N$ đối với một encoding/heuristic cố định. Không nên diễn giải
hiện tượng này thành bảo đảm rằng Product luôn tốt ở kích thước lớn.

\begin{table*}[htbp]
  \centering
  \caption{Thời gian trung vị SAT theo giây; trong ngoặc là số run thành
  công/tổng số run. TO: timeout; KILL: SIGKILL; ``--'': không chạy.}
  \label{tab:sat-time}
  \resizebox{\textwidth}{!}{\input{\ReportPrefix generated/sat_performance_table.tex}}
\end{table*}

OR-Tools hoàn thành 2/3 lần ở $N=256$ với thời gian thành công khoảng
40--43 giây, sau đó timeout hoàn toàn tại $N=512,1000$. CPLEX CP hoàn thành
3/3 lần tại $N=8,16,32$; CPLEX MIP gặp giới hạn license ở $N=32$. Trong dữ
liệu gốc, Gurobi bị giới hạn tại $N=64,128$; các kết quả đó đã được thay bằng
run Academic WLS hợp lệ. Giới hạn cấp phép CPLEX còn lại phản ánh môi trường,
không phải thất bại thuật toán.

Ba run Commander tại $N=1000$ kết thúc với exit code $-9$, tức SIGKILL. Điều
này phù hợp với tình huống process bị hệ thống Colab dừng do áp lực tài nguyên,
nhưng không có log kernel để khẳng định chắc chắn nguyên nhân là OOM; do đó dữ
liệu được gắn nhãn thận trọng là \texttt{KILLED}.

\subsection{Quy mô SAT encoding}

Hình~\ref{fig:variables} và Hình~\ref{fig:clauses} cho thấy Pairwise không dùng
biến phụ nhưng số clause tăng nhanh. Product tạo nhiều biến hơn nhưng có số
clause thấp hơn rõ rệt ở các kích thước lớn quan sát được. Tại $N=1000$,
Product có 1,298,508 biến và 11,942,164 clause, vẫn hoàn thành với trung vị
28.127 giây. Các run timeout/killed không trả lại thống kê encoding, vì vậy
đồ thị dừng đường tại các điểm thiếu thay vì nội suy qua chúng.

\begin{figure*}[htbp]
  \centering
  \includegraphics[width=\textwidth]{plot_variables.png}
  \caption{Số biến Boolean, bao gồm biến phụ, của các SAT encoding.}
  \label{fig:variables}
\end{figure*}

\begin{figure*}[htbp]
  \centering
  \includegraphics[width=\textwidth]{plot_clauses.png}
  \caption{Số clause CNF của các SAT encoding.}
  \label{fig:clauses}
\end{figure*}

\section{Hạn chế và nguy cơ ảnh hưởng tính hợp lệ}

Thứ nhất, Colab là môi trường chia sẻ; mô hình CPU và tổng RAM chưa được lưu
trong bộ metadata cũ, còn Gurobi WLS chạy trong một phiên sau đó. Không nên
diễn giải thứ hạng thời gian như một benchmark được kiểm soát phần cứng. Thứ
hai, số lần lặp khác nhau giữa các phase (1, 3 hoặc 5), và SAT không dùng seed
ngẫu nhiên. Thứ ba, timeout ở process cha làm mất số biến/clause và không cho
biết run dừng trong encoding hay solving. Thứ tư, giới hạn license CPLEX MIP
và RAM Colab làm thu hẹp miền kích thước có thể so sánh; Gurobi chỉ được
chạy bổ sung đến $N=256$. Cuối cùng, nghiên cứu chỉ tìm một lời
giải khả thi; kết quả không đại diện cho bài toán đếm toàn bộ nghiệm N-Queens.

\section{Kết luận}

Project đã triển khai đầy đủ các nhóm phương pháp được yêu cầu và xác minh tính
đúng đắn trên các trường hợp SAT/UNSAT nhỏ. Bộ dữ liệu 241 run đến $N=1000$
cho thấy
việc chọn encoding có ảnh hưởng quyết định: Pairwise có ít biến nhưng bùng nổ
clause, Sequential suy giảm mạnh từ $N=64$, trong khi Product và Binary mở rộng
tốt hơn ở một số kích thước lớn. Gurobi Academic WLS giải thành công toàn bộ
run bổ sung tại $N=64,128,256$; song do thực nghiệm qua nhiều phiên Colab,
không thể suy rộng thứ hạng thời gian thành ưu thế thuật toán tuyệt đối. Các
bất thường không đơn điệu nhấn mạnh rằng không thể đánh giá encoding chỉ bằng
số biến hoặc số clause; cần báo cáo đồng thời thời gian, tỷ lệ hoàn thành và
giới hạn môi trường.

\section*{Data availability}

Mã nguồn, unit test, hướng dẫn Colab, dữ liệu thô, bảng tổng hợp và biểu đồ được
lưu tại \url{\RepositoryURL}. File \texttt{benchmark\_raw.csv} được xem là dữ
liệu gốc; thư mục \texttt{combined\_licensed} chứa bản gộp dùng trong báo cáo
này, còn các thư mục phase giữ nguyên dữ liệu trước khi thay thế run bị
giới hạn license. Mọi bảng và hình trong báo cáo có thể tái tạo bằng
\texttt{combine\_results.py}, \texttt{plot\_results.py} và
\texttt{prepare\_report.py}.

\section*{CRediT authorship contribution statement}

Bùi Đình Cảnh: Conceptualization, Methodology, Software, Validation,
Investigation, Data curation, Visualization, Writing -- original draft,
Writing -- review and editing.

\section*{Declaration of competing interest}

Tác giả tuyên bố không có lợi ích tài chính hoặc quan hệ cá nhân có thể
ảnh hưởng tới công trình này.

\section*{Tuyên bố về công cụ AI hỗ trợ}

Trong quá trình chuẩn bị project, tác giả sử dụng ChatGPT/Codex để hỗ trợ rà
soát mã nguồn, tổ chức thí nghiệm, trực quan hóa và hiệu đính ngôn ngữ. Tác giả
đã kiểm tra mã, nghiệm, số liệu và chịu trách nhiệm cuối cùng về nội dung báo
cáo.

\bibliographystyle{cas-model2-names}
\bibliography{\ReportPrefix references}

\end{document}
